\documentclass[12pt,a4paper]{article}
\usepackage[UTF8]{ctex}
\usepackage{amsmath,amssymb}
\usepackage{geometry}
\usepackage{hyperref}
\usepackage{listings}
\usepackage{xcolor}
\usepackage{booktabs}

\geometry{left=2.5cm,right=2.5cm,top=2.5cm,bottom=2.5cm}

\hypersetup{colorlinks=true,linkcolor=blue,urlcolor=blue}

\lstset{
    language=Python,
    basicstyle=\ttfamily\small,
    keywordstyle=\color{blue},
    commentstyle=\color{green!60!black},
    breaklines=true,
    frame=single,
    backgroundcolor=\color{gray!5},
    showstringspaces=false,
}

\definecolor{certcolor}{RGB}{0,100,0}
\definecolor{uncertcolor}{RGB}{180,0,0}
\newcommand{\certain}[1]{{\color{certcolor}\textbf{[已确认]}} #1}
\newcommand{\uncertain}[1]{{\color{uncertcolor}\textbf{[待验证]}} #1}

\title{\textbf{hw\_pipeline 技术报告}\\广义相对论题目泛化流水线}
\author{}
\date{2026年5月}

\begin{document}
\maketitle
\tableofcontents
\newpage

%=============================================================================
\section{概述}
%=============================================================================

hw\_pipeline 从一道种子 GR 题目出发，自动生成多道可独立求解的新题目。核心思路是两阶段泛化：Phase 1 保持原度规、LLM 自选方向延伸种子题；Phase 2 将题目适配到新度规。正确性由 SymPy/EinsteinPy 的 Fact Sheet 保证。

\certain{流水线架构已实现并可运行。} \uncertain{题目质量尚未经人工评审。}

%=============================================================================
\section{架构}
%=============================================================================

五个子命令：

\begin{lstlisting}
python main.py run input.md --num 3 --subs 2   # 一步完成
python main.py abstract input.md                 # 自然语言 → JSON
python main.py generate seed.json                # 泛化生成
python main.py validate variants/*.json          # 验证+修正
python main.py clean variants/*.json             # 清理排版
\end{lstlisting}

数据流：

\begin{center}
\begin{tabular}{lll}
\toprule
阶段 & 输入 & 输出 \\
\midrule
abstract & PDF/Image/MD & 种子 JSON \\
缓存 & 种子 JSON & 种子 JSON + cached\_fact\_sheet \\
Phase 1 & 种子 JSON + Fact Sheet & $N$ 道基础题 JSON \\
Phase 2 & 基础题 + 新度规 Fact Sheet & $N \times M$ 道替换题 JSON \\
validate & 题目 JSON & validated JSON \\
\bottomrule
\end{tabular}
\end{center}

%=============================================================================
\section{两阶段泛化}
%=============================================================================

\subsection{Phase 1：同度规延伸}

\certain{保持种子题的度规不变，LLM 自选泛化方向编题：}

\begin{itemize}
    \item LLM（COMPOSE prompt）基于 Fact Sheet + 种子题自行选择有意义的延伸方向
    \item 强制 physical\_data 使用种子题的度规（覆写 LLM 输出）
    \item 输出必须是独立完整题目（含物理背景、度规定义、所有假设）
    \item 每道题只能有一问，难度与种子题相当
\end{itemize}

\subsection{Phase 2：度规替换}

\certain{将基础题适配到新度规，保持题目类型和结构不变：}

\begin{itemize}
    \item LLM（PICK\_METRIC）一次性从度规库挑选 $M$ 个替换度规
    \item 所有基础题共享同一组替换度规（减少 API 调用）
    \item LLM（SUBSTITUTE）适配原题到新度规
    \item 替换后 Fact Sheet 由 SymPy 重新计算
\end{itemize}

\uncertain{度规替换的物理合理性依赖 LLM 适配能力。不同物理场景（黑洞 vs 宇宙学）的转换可能不自然。}

\subsection{度规库}

8 个核心度规（均非 heavy，Fact Sheet $<30$s）：

\begin{center}
\begin{tabular}{ll}
\toprule
度规 & 描述 \\
\midrule
Schwarzschild & 静态球对称黑洞，真空解 \\
Minkowski / MinkowskiCartesian / MinkowskiPolar & 平直时空三种坐标 \\
DeSitter & 正宇宙学常数 \\
AntiDeSitter / AntiDeSitterStatic & 负宇宙学常数 \\
FLRW & 标度因子 $a(t)$ \\
\bottomrule
\end{tabular}
\end{center}

Heavy 度规（Kerr、Ernst、KerrNewman、AlcubierreWarp、BesselGravitationalWave、ReissnerNordstrom）排除在替换候选外。

\uncertain{FLRW 含 $a(t)$ 函数，Fact Sheet 分量可能含其导数，LLM 编题时能否正确使用尚不确定。}

%=============================================================================
\section{Fact Sheet}
%=============================================================================

EinsteinPy 符号计算生成：Christoffel $\Gamma^k_{ij}$、Riemann $R^l_{ijk}$、Ricci $R_{ij}/R$、Kretschmann $K=R_{ijkl}R^{ijkl}$、Einstein $G_{\mu\nu}$。所有分量经 sympy.simplify 简化后存为 LaTeX。

\certain{符号计算确定性，给定度规和变量，结果唯一且绝对正确。}

\certain{Fact Sheet 可缓存到种子 JSON 的 cached\_fact\_sheet 字段。run 子命令自动缓存。}

\uncertain{LaTeX→SymPy 转换（用于无 metric\_sympy 的旧 JSON）对非标准格式的鲁棒性有限。流水线尽量使用 metric\_sympy 或度规库直接提供的 SymPy 格式。}

%=============================================================================
\section{验证-修正循环}
%=============================================================================

\subsection{验证}

VALIDATE prompt 聚焦\textbf{可解性和自洽性}：
\begin{enumerate}
    \item 可解性：目标能否推导？缺少必要条件？
    \item 自洽性：度规与题干一致？中间量与 Fact Sheet 一致？
    \item 完整性：做题者能否独立完成？
\end{enumerate}

\certain{answer 为"无答案"→ verified=false（不可解）。}

\subsection{修正循环}

validate\_and\_fix\_loop（max\_attempts=2）：

\begin{enumerate}
    \item 计算 Fact Sheet（只算一次）
    \item validate → 若通过则结束
    \item 若失败且有 issues → fix → 重新验证
    \item 重复直到通过或达到上限
\end{enumerate}

修正规则：度规和 metadata.id/source 不可修改；只修正 origin（answer/solution/hint）；严重错误→ answer="无答案"。

\certain{三个退出路径均添加 tools\_used 标记：EinsteinPy + validate:\{model\}。}

\uncertain{修正循环依赖 LLM 修正能力，可能引入新错误。LLM 验证的判断标准不明确，不同调用可能给出不同结论。}

%=============================================================================
\section{LLM 配置}
%=============================================================================

各阶段模型通过 .env 环境变量配置，未设置时回退到 DEFAULT\_MODEL：

\begin{center}
\begin{tabular}{ll}
\toprule
阶段 & 当前模型 \\
\midrule
abstract / pick\_metric & qwen3.6-flash \\
compose / substitute / validate / fix & glm-5.1 \\
\bottomrule
\end{tabular}
\end{center}

API：DashScope Anthropic 兼容接口，指数退避重试（最多 3 次），超时 600s，max\_tokens 32768。

\certain{JSON 解析鲁棒性：剔除 markdown 代码块 + 修正无效 LaTeX 转义 + 字段过滤。}

\uncertain{API 并发限制未知；不同模型的 JSON 输出格式稳定性待验证。}

%=============================================================================
\section{JSON Schema}
%=============================================================================

\begin{lstlisting}
{
  "metadata": {
    "id": "test-gen-0",
    "type": "calculate",        // prove | calculate | concept
    "tags": { metric, target_object, coordinate, scenario, method },
    "tools_used": ["EinsteinPy", "validate:glm-5.1"],
    "validated": true,
    "source": "generated"
  },
  "physical_data": {
    "dimension": 4,
    "variables": ["$t$", "$r$", "$\\theta$", "$\\phi$"],
    "target": ["$\\Gamma_{01}^{1}$"],
    "metric": [[...], ...],            // LaTeX NxN
    "metric_sympy": [[...], ...]       // SymPy NxN
  },
  "origin": {
    "question": "...",     // 独立完整题干
    "answer": "...",       // 结论或"无答案"
    "solution": "...",     // 推导思路或 null
    "hint": ["..."]        // 提示列表或 null
  }
}
\end{lstlisting}

%=============================================================================
\section{Prompts 关键规则}
%=============================================================================

\begin{center}
\begin{tabular}{lp{8cm}}
\toprule
阶段 & 关键规则 \\
\midrule
ABSTRACT & origin 保留原文照抄；每题一问 \\
COMPOSE & LLM 自选泛化方向；度规不可修改；独立完整题；难度与种子题相当 \\
SUBSTITUTE & 保持题目类型和结构；用新度规数据替换 \\
PICK\_METRIC & 不选当前度规和 heavy 度规 \\
VALIDATE & 聚焦可解性+自洽性；不可解→false \\
FIX & 度规和 metadata 不可修改；严重错误→"无答案" \\
\bottomrule
\end{tabular}
\end{center}

所有 prompt 限制题目范围：只涉及 GR 张量计算与基本应用，严禁引力波、量子引力、宇宙学精细模型等。

solution 引用几何量直接写公式（不标注来源），hint 只写符号约定和特殊简化。

\uncertain{Prompt 约束无法完全防止 LLM 偏离要求，代码层覆写是最后一道防线。}

%=============================================================================
\section{Cleaner}
%=============================================================================

两步清理：1) \_remove\_newlines — 递归移除所有字符串中的 \lstinline|\n|；2) \_clean\_string — 只在 \lstinline|$$...$$| 和 \lstinline|$...$| 数学段内移除空格。metric\_sympy 不参与第二步。

\uncertain{嵌套数学段的正则匹配可能有潜在风险。}

%=============================================================================
\section{正确性保证}
%=============================================================================

三层机制：

\begin{enumerate}
    \item \textbf{SymPy Fact Sheet}：确定性符号计算，绝对正确
    \item \textbf{LLM 验证-修正}：Fact Sheet 为基准，LLM 检查一致性
    \item \textbf{代码覆写}：physical\_data 从种子/度规库获取，不依赖 LLM
\end{enumerate}

\uncertain{三层机制保证度规正确性和结构完整性，但不保证 answer/solution 的物理正确性。LLM 可能引用正确数据得出错误结论。}

%=============================================================================
\section{项目结构}
%=============================================================================

\begin{lstlisting}[basicstyle=\ttfamily\footnotesize]
hw_pipeline/
  main.py            # CLI 入口
  .env / .env.example
  schema/schema.py   # Problem dataclass
  core/
    config.py        # 环境变量配置
    api_client.py    # DashScope API (Anthropic 协议)
    formatter.py     # abstract: 自然语言 → JSON
    pipeline.py      # generate: 两阶段泛化
    system_prompts.py # 6 个 LLM prompts
    sympy_engine.py  # EinsteinPy Fact Sheet
    metric_library.py # 度规库
    cleaner.py       # JSON 清理
    validator.py     # 验证-修正循环
    file_converter.py # PDF/Word → 图片
\end{lstlisting}

%=============================================================================
\section{已解决问题}
%=============================================================================

\begin{enumerate}
    \item LLM 返回额外字段 → dataclass 字段过滤
    \item LaTeX 转义错误 → 双重转义规则 + JSON 解析 fallback
    \item Fact Sheet 重复计算 → 缓存 cached\_fact\_sheet
    \item 3 维度规 Fact Sheet → 限制度规库为 4 维
    \item ReissnerNordstrom eps\_0 → 标记 heavy 并排除
    \item 速度瓶颈 → pick\_metric 一次性调用；并发执行；Phase 1 不生成新度规
\end{enumerate}

%=============================================================================
\section{待验证问题}
%=============================================================================

\begin{enumerate}
    \item \uncertain{LLM 验证可靠性：能否准确判断可解性？需要人工抽样验证准确率。}
    \item \uncertain{修正循环有效性：修正后是否比修正前更好？还是引入新错误？}
    \item \uncertain{度规替换的物理合理性：不同场景转换是否自然？}
    \item \uncertain{难度一致性：LLM 能否准确把握难度？}
    \item \uncertain{FLRW 的 $a(t)$ 导数：LLM 能否正确使用含导数的 Fact Sheet？}
    \item \uncertain{API 并发限制：DashScope rate limit 未知。}
    \item \uncertain{LLM 自选方向的质量：去掉方向池后，LLM 选择的泛化方向是否有意义且多样化？是否会趋向某些固定方向？}
\end{enumerate}

%=============================================================================
\section{后续方向}
%=============================================================================

\begin{enumerate}
    \item \certain{人工评审基准：统计 validate 准确率和题目质量分布。}
    \item \uncertain{heavy 度规支持：优化 Kretschmann 计算，使 Kerr 可用于替换。}
    \item \uncertain{structured output：用 JSON schema 约束 LLM 输出格式。}
    \item \uncertain{多模型交叉验证：减少单模型偏差。}
    \item \uncertain{题目质量评分：教学价值、难度一致性、题干清晰度。}
\end{enumerate}

%=============================================================================
\section*{附注：确定性标注}
\addcontentsline{toc}{section}{附注：确定性标注}
%=============================================================================

{\color{certcolor}\textbf{[已确认]}} — 代码验证或逻辑推理确认，可信度高。
{\color{uncertcolor}\textbf{[待验证]}} — 推测或有限测试，需进一步验证，是后续审阅的重点。

\end{document}